\documentclass[10pt,a4paper]{amsart}
\usepackage{amsmath,amssymb,amsthm,graphicx}
\usepackage{verbatim}
\newenvironment{reference}{\comment}{\endcomment}
%\newenvironment{reference}{}{}
\newenvironment{slogan}{\comment}{\endcomment}
\newenvironment{history}{\comment}{\endcomment}

% For commutative diagrams we use Xy-pic
\usepackage[all]{xy}

% We use 2cell for 2-commutative diagrams.
\xyoption{2cell}
% Theorem environments.
%
\theoremstyle{plain}
\newtheorem{theorem}[subsection]{Theorem}
\newtheorem{proposition}[subsection]{Proposition}
\newtheorem{lemma}[subsection]{Lemma}

\theoremstyle{definition}
\newtheorem{definition}[subsection]{Definition}
\newtheorem{example}[subsection]{Example}
\newtheorem{exercise}[subsection]{Exercise}
\newtheorem{situation}[subsection]{Situation}

\theoremstyle{remark}
\newtheorem{remark}[subsection]{Remark}
\newtheorem{remarks}[subsection]{Remarks}

\numberwithin{equation}{subsection}

% Macros
%
\def\lim{\mathop{\mathrm{lim}}\nolimits}
\def\colim{\mathop{\mathrm{colim}}\nolimits}
\def\Spec{\mathop{\mathrm{Spec}}}
\def\Hom{\mathop{\mathrm{Hom}}\nolimits}
\def\Ext{\mathop{\mathrm{Ext}}\nolimits}
\def\SheafHom{\mathop{\mathcal{H}\!\mathit{om}}\nolimits}
\def\SheafExt{\mathop{\mathcal{E}\!\mathit{xt}}\nolimits}
\def\Sch{\mathit{Sch}}
\def\Mor{\mathop{\mathrm{Mor}}\nolimits}
\def\Ob{\mathop{\mathrm{Ob}}\nolimits}
\def\Sh{\mathop{\mathit{Sh}}\nolimits}
\def\NL{\mathop{N\!L}\nolimits}
\def\CH{\mathop{\mathrm{CH}}\nolimits}
\def\proetale{{pro\text{-}\acute{e}tale}}
\def\etale{{\acute{e}tale}}
\def\QCoh{\mathit{QCoh}}
\def\Ker{\mathop{\mathrm{Ker}}}
\def\Im{\mathop{\mathrm{Im}}}
\def\Coker{\mathop{\mathrm{Coker}}}
\def\Coim{\mathop{\mathrm{Coim}}}

% Boxtimes
%
\DeclareMathSymbol{\boxtimes}{\mathbin}{AMSa}{"02}

%
% Macros for moduli stacks/spaces
%
\def\QCohstack{\mathcal{QC}\!\mathit{oh}}
\def\Cohstack{\mathcal{C}\!\mathit{oh}}
\def\Spacesstack{\mathcal{S}\!\mathit{paces}}
\def\Quotfunctor{\mathrm{Quot}}
\def\Hilbfunctor{\mathrm{Hilb}}
\def\Curvesstack{\mathcal{C}\!\mathit{urves}}
\def\Polarizedstack{\mathcal{P}\!\mathit{olarized}}
\def\Complexesstack{\mathcal{C}\!\mathit{omplexes}}
% \Pic is the operator that assigns to X its picard group, usage \Pic(X)
% \Picardstack_{X/B} denotes the Picard stack of X over B
% \Picardfunctor_{X/B} denotes the Picard functor of X over B
\def\Pic{\mathop{\mathrm{Pic}}\nolimits}
\def\Picardstack{\mathcal{P}\!\mathit{ic}}
\def\Picardfunctor{\mathrm{Pic}}
\def\Deformationcategory{\mathcal{D}\!\mathit{ef}}
\usepackage{cleveref}
\crefname{originallemma}{original lemma}{original lemmas}
\Crefname{originallemma}{Original lemma}{Original lemmas}
\makeatletter
\newlabel{lemma-going-up-closed}{{00HZ}{source}{}{sourcepointer-1}{}}
\newlabel{lemma-going-up-closed@cref}{{[originallemma][0][]00HZ}{[1][0][]source}{}{}{}}
\makeatother
\crefname{originallemma}{original lemma}{original lemmas}
\Crefname{originallemma}{Original lemma}{Original lemmas}
\makeatletter
\newlabel{lemma-image-dense-generic-points}{{00FL}{source}{}{sourcepointer-2}{}}
\newlabel{lemma-image-dense-generic-points@cref}{{[originallemma][0][]00FL}{[1][0][]source}{}{}{}}
\makeatother
\crefname{originallemma}{original lemma}{original lemmas}
\Crefname{originallemma}{Original lemma}{Original lemmas}
\makeatletter
\newlabel{lemma-integral-going-up}{{00GU}{source}{}{sourcepointer-3}{}}
\newlabel{lemma-integral-going-up@cref}{{[originallemma][0][]00GU}{[1][0][]source}{}{}{}}
\makeatother
\crefname{originallemma}{original lemma}{original lemmas}
\Crefname{originallemma}{Original lemma}{Original lemmas}
\makeatletter
\newlabel{lemma-surjective-spec-radical-ideal}{{00FI}{source}{}{sourcepointer-4}{}}
\newlabel{lemma-surjective-spec-radical-ideal@cref}{{[originallemma][0][]00FI}{[1][0][]source}{}{}{}}
\makeatother
\title{Generators whose square and cube lie in the base, with locally nilpotent kernel, give a universal homeomorphism}
\author{The Stacks Project Authors}
\date{stacks-project:algebra.tex:local-snapshot-20261012; GNU FDL 1.2 or later}
\begin{document}
\maketitle
\noindent\textit{Editorial note.} All delivered mathematical text is the author's own, copied from the pinned Stacks Project chapter file algebra.tex; no statement or proof step was rewritten, reordered or removed.\par\medskip
\noindent\textit{Editorial note.} Preamble provenance, as required for this source: the frozen snapshot stages the chapter's own preamble file beside it (preamble.tex, 217 lines, 8116 bytes, the same shared Stacks-wide preamble the published chow and etale-cohomology cards use); the first line of the pinned chapter file is \textbackslash{}input\{preamble\}, so no preamble text lives inside the chapter itself. The card takes that file's lines 8-18 (the verbatim package and the three \textbackslash{}comment-based environments reference/slogan/history at lines 9-12, section 13.2, together with the Xy-pic setup) and its lines 154 to the runtime end of the file (the native theorem declarations and mathematical macros). Not one character was changed, no package was added, no macro was defined, and the excerpt serves only as preamble, never as a span coordinate.\par\medskip
\noindent\textit{Editorial note.} The same-chapter pointers are located at their \textbackslash{}label lines inside the pinned chapter file; the cross-chapter pointers are located in the project's own tag list staged in the snapshot (sources/stacks-chow/tags.txt), at the verbatim lines carrying their canonical tags.\par\medskip
\noindent\textit{Editorial note.} The counted proof refers to same-chapter statements by label; they are recorded as navigation stubs pointed at their label lines.\par\medskip
\noindent\textit{Editorial note.} The counted spans contain no cite, so the card carries no bibliography.\par\medskip
\noindent\textit{Allowed standard prerequisites (editorial).}\begin{itemize}
\item Closedness of the map on spectra for integral ring maps. Invoked in the counted proof at lines 11000-11002. Gives closedness of the induced map on spectra.
\item Density of the image from generic points. Invoked in the counted proof at lines 11002-11003 and again at the end. Supplies density of the image, which together with closedness gives surjectivity.
\end{itemize}\medskip
\ifdefined\hypertarget\hypertarget{sourcepointer-1}{}\fi
\noindent\textit{Source reference.} Original lemma 00HZ: sources/stacks-algebra/pinned-original-source.tex, source lines 9869--9869. Navigation pointer for \textbackslash{}ref\{lemma-going-up-closed\} used inside the counted span (line 11001); the argument that consumes it is retained in the span.\par\medskip
\ifdefined\hypertarget\hypertarget{sourcepointer-2}{}\fi
\noindent\textit{Source reference.} Original lemma 00FL: sources/stacks-algebra/pinned-original-source.tex, source lines 6138--6138. Navigation pointer for \textbackslash{}ref\{lemma-image-dense-generic-points\} used inside the counted span (line 11003); the argument that consumes it is retained in the span.\par\medskip
\ifdefined\hypertarget\hypertarget{sourcepointer-3}{}\fi
\noindent\textit{Source reference.} Original lemma 00GU: sources/stacks-algebra/pinned-original-source.tex, source lines 8070--8070. Navigation pointer for \textbackslash{}ref\{lemma-integral-going-up\} used inside the counted span (line 11002); the argument that consumes it is retained in the span.\par\medskip
\ifdefined\hypertarget\hypertarget{sourcepointer-4}{}\fi
\noindent\textit{Source reference.} Original lemma 00FI: sources/stacks-algebra/pinned-original-source.tex, source lines 6039--6039. Navigation pointer for \textbackslash{}ref\{lemma-surjective-spec-radical-ideal\} used inside the counted span (line 11022); the argument that consumes it is retained in the span.\par\medskip
\begin{lemma}
\label{lemma-2-3-ring-map}
Let $\varphi : R \to S$ be a ring map. Assume
\begin{enumerate}
\item[(a)] $S$ is generated as an $R$-algebra by elements $x$ such
that $x^2, x^3 \in \varphi(R)$, and
\item[(b)] $\Ker(\varphi)$ is locally nilpotent,
\end{enumerate}
Then $\varphi$ induces isomorphisms on residue fields and
a homeomorphism of spectra. For any ring map $R \to R'$
the ring map $R' \to R' \otimes_R S$ also satisfies (a) and (b).
\end{lemma}

\begin{proof}
Assume (a) and (b). The map on spectra is closed as $S$ is integral
over $R$, see Lemmas \ref{lemma-going-up-closed} and
\ref{lemma-integral-going-up}. The image is dense by
Lemma \ref{lemma-image-dense-generic-points}. Thus $\Spec(S) \to \Spec(R)$
is surjective. If $\mathfrak q \subset S$ is a prime lying over
$\mathfrak p \subset R$ then the field extension
$\kappa(\mathfrak q)/\kappa(\mathfrak p)$ is generated by elements
$\alpha \in \kappa(\mathfrak q)$ whose square and cube are
in $\kappa(\mathfrak p)$. Thus clearly $\alpha \in \kappa(\mathfrak p)$
and we find that $\kappa(\mathfrak q) = \kappa(\mathfrak p)$.
If $\mathfrak q, \mathfrak q'$ were two distinct primes lying over
$\mathfrak p$, then at least one of the generators $x$ of $S$ as
in (a) would have distinct images in
$\kappa(\mathfrak q) = \kappa(\mathfrak p)$ and
$\kappa(\mathfrak q') = \kappa(\mathfrak p)$.
This would contradict the fact that both $x^2$ and $x^3$
do have the same image. This proves that $\Spec(S) \to \Spec(R)$
is injective hence a homeomorphism (by what was already shown).

\medskip\noindent
Since $\varphi$ induces a homeomorphism on spectra, it is in particular
surjective on spectra which is a property preserved under any base change, see
Lemma \ref{lemma-surjective-spec-radical-ideal}.
Therefore for any $R \to R'$ the kernel of the ring map
$R' \to R' \otimes_R S$ consists of nilpotent elements, see
Lemma \ref{lemma-image-dense-generic-points},
in other words (b) holds for $R' \to R' \otimes_R S$.
It is clear that (a) is preserved under base change.
\end{proof}

\end{document}
